% ============================================================
%  05_kiem_thu_hop_trang.tex  --  Danh gia rui ro (CVSS-lite)
% ============================================================

\section{Đánh giá rủi ro}
\label{sec:hoptrang}

Kiểm thử hộp trắng ở đây bao gồm việc đọc mã nguồn, hiểu cơ chế lỗi, và từ đó
đánh giá định lượng mức độ rủi ro của từng lỗ hổng. Script
\textit{python/risk\_assessor.py} thực hiện tính toán tự động trên dataset.

% ────────────────────────────────────────────────────────────
\subsection{Mô hình CVSS-lite được áp dụng}

Script tính hai chỉ số phụ rồi tổng hợp thành điểm cơ sở:

\begin{align}
    \mathrm{Exploitability}_{sub} &= 8.22 \times AV \times AC \times PR \times UI \\
    \mathrm{ISC}_{base} &= 1 - (1-C)(1-I)(1-A) \\
    \mathrm{Impact}_{sub} &= 6.42 \times \mathrm{ISC}_{base} \\
    \mathrm{BaseScore} &= \min(\mathrm{Impact}_{sub} + \mathrm{Exploitability}_{sub},\ 10.0)
\end{align}

Trong đó các trọng số được ánh xạ từ chuỗi định tính sang giá trị số:

\begin{table}[H]
\centering
\begin{tabularx}{\textwidth}{|p{3.5cm}|X|}
\hline
\textbf{Thuộc tính} & \textbf{Giá trị trọng số} \\ \hline
Attack Vector (AV)  & NETWORK=0.85, ADJACENT=0.62, LOCAL=0.55, PHYSICAL=0.20 \\ \hline
Attack Complexity (AC) & LOW=0.77, HIGH=0.44 \\ \hline
Privileges Required (PR) & NONE=0.85, LOW=0.62, HIGH=0.27 \\ \hline
User Interaction (UI) & NONE=0.85, REQUIRED=0.62 \\ \hline
C/I/A Impact & NONE=0.00, LOW=0.22, HIGH=0.56 \\ \hline
\end{tabularx}
\caption{Bảng trọng số CVSS-lite}
\end{table}

% ────────────────────────────────────────────────────────────
\subsection{Phân tích từng lỗ hổng}

\subsubsection{VULN-001 -- CWE-121: Buffer Overflow (9.8 / CRITICAL)}

\begin{itemize}
    \item Attack Vector: NETWORK -- lỗi này xảy ra khi dữ liệu đầu vào
          dài được truyền qua mạng (API, form nhập liệu). Không cần truy cập vật
          lý hay cục bộ.
    \item Attack Complexity: LOW -- không cần điều kiện đặc biệt; chỉ
          cần gửi chuỗi dài hơn 31 ký tự.
    \item Privileges Required: NONE -- người dùng ẩn danh cũng khai thác
          được.
    \item User Interaction: NONE -- không cần nạn nhân thực hiện thao
          tác nào.
    \item C/I/A: HIGH/HIGH/HIGH -- buffer overflow có thể dẫn đến thực
          thi mã tùy ý, toàn quyền đọc/ghi/phá hủy hệ thống.
\end{itemize}

Tính toán:
\[
\mathrm{Exp}_{sub} = 8.22 \times 0.85 \times 0.77 \times 0.85 \times 0.85 = 3.89
\]
\[
\mathrm{ISC} = 1 - (1-0.56)^3 = 1 - 0.085 = 0.915 \quad \Rightarrow \quad
\mathrm{Imp}_{sub} = 6.42 \times 0.915 = 5.87
\]
\[
\mathrm{BaseScore} = \min(5.87 + 3.89,\ 10) = \min(9.76,\ 10) \approx 9.8
\]

\subsubsection{VULN-002 -- CWE-416: Use-After-Free (6.9 / MEDIUM)}

\begin{itemize}
    \item Attack Vector: LOCAL -- kẻ tấn công cần khả năng chạy code
          trên cùng máy (heap grooming đòi hỏi kiểm soát luồng cấp phát).
    \item Attack Complexity: HIGH -- cần kỹ thuật heap manipulation
          phức tạp để kiểm soát nội dung vùng nhớ tái cấp phát.
    \item Privileges Required: LOW -- cần tài khoản thường (không root).
    \item User Interaction: NONE.
    \item C/I/A: HIGH/HIGH/HIGH -- nếu khai thác thành công vẫn dẫn
          đến thực thi mã tùy ý.
\end{itemize}

Tính toán:
\[
\mathrm{Exp}_{sub} = 8.22 \times 0.55 \times 0.44 \times 0.62 \times 0.85 = 1.05
\]
\[
\mathrm{BaseScore} = \min(5.87 + 1.05,\ 10) = 6.9
\]

Điểm Impact cao (vì C/I/A đều HIGH) nhưng Exploitability thấp do cần điều kiện
khai thác phức tạp và quyền LOCAL, kéo điểm tổng xuống mức MEDIUM.

\subsubsection{VULN-003 -- CWE-401: Memory Leak (5.4 / MEDIUM)}

\begin{itemize}
    \item Attack Vector: LOCAL, AC: LOW, PR: LOW, UI: NONE.
    \item C/I/A: NONE/NONE/HIGH -- rò rỉ bộ nhớ không trực tiếp làm
          lộ dữ liệu hay phá vỡ tính toàn vẹn; tác động duy nhất là cạn kiệt
          tài nguyên (Availability).
\end{itemize}

Tính toán:
\[
\mathrm{Exp}_{sub} = 8.22 \times 0.55 \times 0.77 \times 0.62 \times 0.85 = 1.83
\]
\[
\mathrm{ISC} = 1 - (1-0)(1-0)(1-0.56) = 0.56 \quad \Rightarrow \quad
\mathrm{Imp}_{sub} = 6.42 \times 0.56 = 3.59 \approx 3.60
\]
\[
\mathrm{BaseScore} = \min(3.60 + 1.83,\ 10) = 5.4
\]

\subsubsection{VULN-004 -- CWE-476: NULL Pointer Dereference (2.5 / LOW)}

\begin{itemize}
    \item Attack Vector: LOCAL, AC: HIGH, PR: LOW, UI: NONE.
    \item C/I/A: NONE/NONE/LOW -- chỉ gây crash (Availability LOW),
          không làm lộ hay sửa dữ liệu.
    \item Đây là lỗi \emph{tiềm ẩn}: phát hiện qua code review chứ chưa được
          ASan bắt (chỉ xảy ra khi \textit{malloc()} thất bại).
\end{itemize}

Tính toán:
\[
\mathrm{Exp}_{sub} = 8.22 \times 0.55 \times 0.44 \times 0.62 \times 0.85 = 1.05
\]
\[
\mathrm{ISC} = 1 - (1)(1)(1-0.22) = 0.22 \quad \Rightarrow \quad
\mathrm{Imp}_{sub} = 6.42 \times 0.22 = 1.41
\]
\[
\mathrm{BaseScore} = \min(1.41 + 1.05,\ 10) = 2.5
\]

\subsubsection{VULN-005 -- CWE-134: Format String Bug (9.8 / CRITICAL)}

Đây là lỗ hổng bổ sung được thêm vào dataset để minh họa một lớp lỗi khác
trong C. Lỗi xảy ra khi chuỗi do người dùng kiểm soát được truyền trực tiếp
làm format string:

\begin{lstlisting}[language=C, caption={Ví dụ lỗi format string -- không bao giờ làm thế này}]
/* LOI: user_input co the chua %x, %s, %n ... */
printf(user_input);         /* KHONG AN TOAN */

/* DUNG: luon dung format string hang so */
printf("%s", user_input);   /* AN TOAN */
\end{lstlisting}

Các chỉ thị \textit{\%x} đọc dữ liệu từ stack (rò rỉ thông tin), \textit{\%s}
có thể đọc tùy ý từ con trỏ trên stack, \textit{\%n} \emph{ghi} số ký tự đã in
vào một địa chỉ bộ nhớ -- đây là primitive đọc/ghi tùy ý mạnh tương đương buffer
overflow.

\begin{itemize}
    \item Attack Vector: NETWORK, AC: LOW, PR: NONE, UI: NONE -- giống VULN-001,
          hoàn toàn từ xa, không đặc quyền.
    \item C/I/A: HIGH/HIGH/HIGH -- có thể đọc bộ nhớ tùy ý, ghi bộ nhớ tùy ý
          qua \textit{\%n}, gây crash.
\end{itemize}

Tính toán: Cùng vector với VULN-001 $\Rightarrow$ BaseScore = 9.8 (CRITICAL).

% ────────────────────────────────────────────────────────────
\subsection{Bảng đánh giá rủi ro tổng hợp}

\begin{table}[H]
\centering
\begin{tabularx}{\textwidth}{|p{1.6cm}|p{1.7cm}|X|p{0.9cm}|p{1.3cm}|p{1.2cm}|p{1.3cm}|}
\hline
\textbf{ID} & \textbf{CWE} & \textbf{Tên lỗ hổng} & \textbf{Exp} & \textbf{Imp} & \textbf{Điểm} & \textbf{Mức độ} \\ \hline
VULN-001 & CWE-121 & Buffer Overflow & 3.89 & 5.87 & 9.8 & CRITICAL \\ \hline
VULN-005 & CWE-134 & Format String Bug & 3.89 & 5.87 & 9.8 & CRITICAL \\ \hline
VULN-002 & CWE-416 & Use-After-Free & 1.05 & 5.87 & 6.9 & MEDIUM \\ \hline
VULN-003 & CWE-401 & Memory Leak & 1.83 & 3.60 & 5.4 & MEDIUM \\ \hline
VULN-004 & CWE-476 & NULL Ptr Deref & 1.05 & 1.41 & 2.5 & LOW \\ \hline
\end{tabularx}
\caption{Kết quả đánh giá rủi ro CVSS-lite -- 5 lỗ hổng}
\end{table}

% ────────────────────────────────────────────────────────────
\subsection{Thứ tự ưu tiên khắc phục}

\[
\text{VULN-001} \rightarrow \text{VULN-005} \rightarrow \text{VULN-002}
\rightarrow \text{VULN-003} \rightarrow \text{VULN-004}
\]

Cách sắp xếp này phù hợp với nguyên tắc bảo mật thực tế:

\begin{enumerate}
    \item VULN-001 (Buffer Overflow) -- Ưu tiên số 1 tuyệt đối. Đây là lỗi thật
          sự hiện diện trong module, có thể bị khai thác từ xa không cần đặc
          quyền, tiềm năng dẫn đến thực thi mã tùy ý. Sửa bằng cách thay
          \textit{strcpy} bằng \textit{snprintf}.
    \item VULN-005 (Format String) -- Cùng điểm 9.8 nhưng là kịch bản mở rộng
          chưa hiện diện trực tiếp trong code hiện tại; cần lưu ý khi phát triển
          module logging hoặc hiển thị tên sách.
    \item VULN-002 (Use-After-Free) -- Điểm Impact cao (5.87) nhưng khai thác
          phức tạp. Sửa bằng cách sao chép dữ liệu trước \textit{free} và đặt
          con trỏ về \textit{NULL}.
    \item VULN-003 (Memory Leak) -- Chỉ ảnh hưởng Availability, không rò rỉ dữ
          liệu nhạy cảm. Tuy nhiên trong hệ thống chạy liên tục, cần sửa sớm để
          tránh cạn tài nguyên.
    \item VULN-004 (NULL Ptr Deref) -- Điểm thấp nhất, chỉ xảy ra khi hệ thống
          hết bộ nhớ (tình huống hiếm). Thêm kiểm tra \textit{if (!buf)} trước
          khi dùng con trỏ là đủ.
\end{enumerate}

% ────────────────────────────────────────────────────────────
\subsection{Biểu đồ điểm rủi ro}

\begin{figure}[H]
\centering
\begin{tikzpicture}
\begin{axis}[
    xbar,
    width=0.85\textwidth,
    height=7.5cm,
    xlabel={Điểm rủi ro BaseScore},
    ytick={1,2,3,4,5},
    yticklabels={VULN-004,VULN-003,VULN-002,VULN-005,VULN-001},
    xmin=0, xmax=11,
    nodes near coords,
    nodes near coords align={horizontal},
    every node near coord/.append style={font=\footnotesize},
    bar width=0.45cm,
    enlarge y limits=0.15,
    tick label style={font=\small},
    label style={font=\small},
]
\addplot[fill=gray!50] coordinates {
    (2.5, 1)
    (5.4, 2)
    (6.9, 3)
    (9.8, 4)
    (9.8, 5)
};
\end{axis}
\end{tikzpicture}
\caption{Biểu đồ điểm rủi ro CVSS-lite của 5 lỗ hổng}
\end{figure}

% ────────────────────────────────────────────────────────────
\subsection{Khuyến nghị khắc phục}

\begin{table}[H]
\centering
\begin{tabularx}{\textwidth}{|p{1.6cm}|X|}
\hline
\textbf{Lỗ hổng} & \textbf{Hành động khắc phục} \\ \hline
VULN-001 & Thay toàn bộ \textit{strcpy} bằng \textit{snprintf(dst, sizeof(dst), "\%s", src)}; kiểm tra giá trị trả về; từ chối đầu vào quá dài thay vì im lặng cắt bớt \\ \hline
VULN-005 & Luôn dùng format string hằng: \textit{printf("\%s", user\_str)} thay vì \textit{printf(user\_str)}; bật cảnh báo \textit{-Wformat-security} trong quá trình biên dịch \\ \hline
VULN-002 & Sao chép dữ liệu cần thiết trước khi \textit{free()}; đặt con trỏ về \textit{NULL} ngay sau \textit{free()}; sử dụng ownership rõ ràng \\ \hline
VULN-003 & Đảm bảo mọi nhánh code (kể cả nhánh lỗi) đều giải phóng bộ nhớ; dùng pattern ``cấp phát -- dùng -- giải phóng'' nhất quán \\ \hline
VULN-004 & Kiểm tra giá trị trả về của \textit{malloc()} trước khi dùng con trỏ; xử lý trường hợp cấp phát thất bại một cách rõ ràng \\ \hline
\end{tabularx}
\caption{Khuyến nghị khắc phục cho 5 lỗ hổng}
\end{table}
